\documentclass[11pt]{amsart}
\usepackage[T1]{fontenc}
\usepackage{lmodern}
\usepackage{amsmath,amssymb,amsthm}
\usepackage{xcolor}
\usepackage[colorlinks=true,linkcolor=blue!55!black,citecolor=blue!55!black,urlcolor=blue!55!black]{hyperref}
\newtheorem{theorem}{Theorem}[section]
\newtheorem{proposition}[theorem]{Proposition}
\newcommand{\set}[1]{\left\{#1\right\}}

\title{The Exact Value of a Pairwise-GCD Sunflower-Free Function at 126}
\author{Zijian Zeng}
\address{Institute of Computer Science and Digital Innovation, UCSI University, Kuala Lumpur, 56000, MALAYSIA}
\email{zijianzeng@foxmail.com; 1002266693@ucsiuniversity.edu.my}
\date{}
\hypersetup{pdftitle={The Exact Value of a Pairwise-GCD Sunflower-Free Function at 126},pdfauthor={Zijian Zeng},pdfsubject={An exact finite result in extremal number theory}}

\begin{document}
\begin{abstract}
For a positive integer $N$, let $f_3(N)$ be the largest cardinality of a set
$A\subseteq\{1,\ldots,N\}$ containing no three distinct elements whose three
pairwise greatest common divisors are equal. We prove the exact finite result
$f_3(126)=16$. The lower bound is an explicit set. For the upper bound, the
forbidden triples and an exact cardinality circuit give a Boolean formula whose
unsatisfiability is certified by a checkable DRAT derivation. This result
concerns only the stated finite value and does not resolve the asymptotic
estimation problem.
\end{abstract}
\maketitle

\section{Introduction}
Erd\H{o}s studied the problem of finding many integers without $t$ members
whose pairwise greatest common divisors are all equal
\cite{Erd_s_1964}. His notation used the least length forcing such a
configuration. We use the current maximum-avoiding convention recorded for
Erd\H{o}s Problem 535 \cite{ErdosProblems535}; the two conventions differ by
one.

The classical literature on asymptotic bounds includes work of Abbott and
Hanson \cite{Abbott_1970}. Our scope is deliberately finite.
Searches of the primary records and public exact-data repositories listed in
the accompanying archive, through 1 September 2026, found no verifiable
certificate for the value below. This is a bounded search report, not an
assertion that no unpublished or differently indexed computation exists.

\section{Definition and result}
For $N\geq1$, write $[N]=\{1,2,\ldots,N\}$ and define
\begin{align*}
 f_3(N)=\max\bigl\{|A|: {}& A\subseteq[N],\text{ and for every }
 \set{a,b,c}\in\tbinom{A}{3},\\[-2mm]
 &\gcd(a,b),\gcd(a,c),\gcd(b,c)\text{ are not all equal}\bigr\}.
\end{align*}

\begin{theorem}\label{thm:main}
$f_3(126)=16$.
\end{theorem}

\section{Exact lower bound}
Let
\begin{equation}\label{eq:witness}
A=\set{14,20,21,24,25,30,36,40,48,63,80,90,96,98,108,125}.
\end{equation}
A direct integer calculation checks all $\binom{16}{3}=560$ triples in $A$.
For every such triple, the three pairwise gcds are not all equal. Thus
$f_3(126)\geq16$. The archive includes a short independent verifier that
parses the displayed set and repeats precisely this calculation.

\section{Certified upper bound}
Introduce a Boolean variable $x_i$ for each $i\in[126]$. For every unordered
triple $\set{a,b,c}$ satisfying
\[
 \gcd(a,b)=\gcd(a,c)=\gcd(b,c),
\]
insert the clause $\neg x_a\lor\neg x_b\lor\neg x_c$. There are 114345 such
clauses.

To impose $\sum_i x_i\geq17$, for $1\leq j\leq\min(i,17)$ introduce
$s_{i,j}$ with the exact meaning that at least $j$ of $x_1,\ldots,x_i$ are
true. The boundary cases and recurrence
\[
s_{i,j}\ \longleftrightarrow\
s_{i-1,j}\lor(s_{i-1,j-1}\land x_i)
\]
are encoded in both directions by Tseitin clauses, and $s_{126,17}$ is
asserted. Consequently, the CNF is satisfiable if and only if an admissible set
of size at least 17 exists.

The deterministic generator produces 2132 variables and 122227 clauses. A SAT
solver emitted a DRAT derivation of unsatisfiability. The derivation was then
replayed by an independently distributed DRAT checker, which reported
\texttt{s VERIFIED}. Its checked core used 112162 input clauses, 1834588
lemmas, and 120875397 resolution steps. A semantic cross-check also compares
the CNF answer with direct enumeration of every subset for each $N\leq12$ and
every cardinality threshold.

\begin{proposition}\label{prop:upper}
No admissible subset of $[126]$ has cardinality at least $17$.
\end{proposition}
\begin{proof}
The generator enumerates every unordered triple of distinct members of
$[126]$ and inserts exactly the corresponding forbidden-triple clause. By
induction on $i$, the two-direction recurrence clauses make $s_{i,j}$ true
exactly when at least $j$ of the first $i$ selection variables are true.
Therefore a satisfying assignment would encode an admissible set of size at
least 17. The checked DRAT derivation proves that this CNF is unsatisfiable.
\end{proof}

Combining Proposition~\ref{prop:upper} with the set in
\eqref{eq:witness} proves Theorem~\ref{thm:main}.

\section{Reproducibility and limitations}
The accompanying computational supplement contains the deterministic CNF generator, direct
witness verifier, small-instance cross-check, DIMACS instance, DRAT
certificate, solver log, and independent checker log. Fast checks and the full
certificate replay are separate commands. All decisive arithmetic is
integer-valued; timeouts and heuristic searches play no role in the proof.
The theorem rests on the explicit witness, source code, and independently
checkable certificate described above.

The theorem determines one finite cell only. It neither proves the supplied
claims about earlier finite values nor advances the asymptotic mother problem.

\bibliographystyle{amsplain}
\bibliography{references}
\end{document}
